\section{Validation: Cartel Adjacency Under Coarsened Observability}
\label{sec:cade}

The validation asks the structurally appropriate question for a
screening statistic computed on award-record data: does the
construct concentrate on populations of firms with a higher prior
probability of operating inside cartel-affected environments
than random selection at comparable participation volume? The
question is not whether the construct identifies cartel members.
Award-record data does not contain the information required for
membership identification, and the screening framing of
\S\ref{sec:emp_estimand} does not require it. We treat the
cobidder population---firms that participated alongside CADE
direct defendants---as the empirical content the screening
statistic is designed to recover, and report against
direct-defendant labels separately to bound the construct's
scope.

\subsection{Validation populations}
\label{sec:cade_objects}

\begin{table}[htbp]
\centering
\caption{Validation populations and object discipline}
\label{tab:cade_populations}
\footnotesize
\begin{adjustbox}{max width=\textwidth,center}
\begin{tabular}{p{4.0cm}p{2.0cm}p{6.5cm}}
\toprule
Population & Size & Definition and role \\
\midrule
All BEC participants & $41{,}444$ & Universe of firms bidding in BEC, $2009$--$2019$. \\
Always-losers & $\valAlwaysLosers$ & Firms with zero wins across the sample; stratum within which the screening statistic operates. \\
Frequent losers & $\valFL$ & Always-losers above the median $+\,1.5\!\times$IQR participation threshold; the flagged population. \\
CADE direct defendants in BEC & $\valDirectCADE$ & Firms named in CADE rulings; typically frequent winners; outside the screening statistic's target population. \\
Cobidders in the always-loser stratum & $\valCobidders$ & Always-losers that participated alongside CADE direct defendants; the cartel-adjacency label the screening statistic is designed to recover. \\
\bottomrule
\end{tabular}
\end{adjustbox}

\smallskip
{\footnotesize\textit{Notes:} The screening statistic operates
within the always-loser stratum, not across the broader BEC
universe. Participation alongside an adjudicated cartel member
is consistent with cover bidding, capacity-transfer motives,
market intelligence, or bystander participation in
cartel-affected markets. Strong AUC against this label
corroborates that the screening statistic concentrates on
cartel-relevant environments; it does not adjudicate
membership.\par}
\end{table}

\subsection{Conservative benchmark}
\label{sec:cade_contemporaneous}

The primary benchmark restricts the ground truth to four CADE
cases adjudicated before $2020$ (coleta de lixo, aquecedores
solares, trens-metr\^os, tecnologia-informa\c{c}\~ao
$2019$-$10$): $30$ firm-defendants, all active in BEC, with $210$
always-losers co-bidding with at least one of them, of which
$108$ are frequent losers. The participation-stratified
co-bidding rate is $\valCADEenrich$ against a baseline of
$\valCADEbase$ ($p<0.001$, $\valPermPilot$ permutations), an
excess ratio of $\valMultThreeTwo\!\!\times$. The ROC AUC
against the contemporaneous ground truth is $\valAUCprePost$
($95\%$ DeLong CI $[0.713, 0.783]$).

This is the primary benchmark for the screening interpretation.
The construct uses no information from cases adjudicated
post-sample and the labels predate estimation: existence and
membership were established by enforcement before the screening
statistic was built. Within the always-loser stratum, frequent
losers co-occur with adjudicated cartel members at rates above
random matching at comparable participation volume. The excess
is what the screening framing predicts under the separating
equilibrium of Online Appendix~A: where the cartel deploys
cover bidders, the loser-side participation footprint
concentrates.

\subsection{Prospective benchmark and within-firm enrichment}
\label{sec:cade_prospective}

A more permissive benchmark uses CADE's full procurement-cartel
portfolio ($12$ cases, $65$ firm-defendants, of which
$\valDirectCADE$ are active in BEC). Among $\valFL$ frequent
losers, $\valCobidders$ ($\valCobidShareFL$) co-participate with
at least one of the $\valDirectCADE$; the
participation-stratified excess ratio is
$\valMultThreeFive\!\!\times$ ($p<0.001$). Firm-level AUC is
$\valAUCFLfirm$ ($95\%$ CI $\valAUCFLfirmCI$) in-sample and
$\valAUCFLfirmTemp$ ($\valAUCFLfirmTempCI$) under temporal
holdout (train $2009$--$2016$, test $2017$--$2019$). The
benchmark treats post-sample adjudications as informative labels
for cartel involvement during the BEC window---a reasonable but
non-trivial assumption that the conservative benchmark in
\S\ref{sec:cade_contemporaneous} avoids. We report it as
complementary, not as the primary anchor.

A within-firm exercise on the same portfolio: $3$ of $7$
always-loser direct defendants ($\valPctFortythree$) cross the
frequent-loser threshold against a population baseline of
$\valPctSixteenRef$. Table~\ref{tab:cade_fl_firms} lists them.

\begin{table}[htbp]
\centering
\caption{Direct CADE defendants flagged by the construct}
\label{tab:cade_fl_firms}
\small
\begin{tabular}{lllcc}
\toprule
Firm & Cartel & CADE process & Tenders & Wins \\
\midrule
Sol Tecnologia & Solar water heaters & 08012.001273/2010-24 & 84 & 0 \\
Nova Esperan\c{c}a Locadora & School transportation & 08700.005876/2019-85 & 97 & 0 \\
Jofran Com\'ercio & Trash bags (Op.\ Colludium) & 08700.005789/2015-02 & 65 & 0 \\
\bottomrule
\end{tabular}
\end{table}

These are the exception within a direct-defendant population
that is otherwise frequent-winner-heavy; the AUC against direct
defendants in the broader BEC firm universe remains
$\valAUCdirectStd$. The asymmetry between cobidder
discrimination ($\valAUCFLfirm$ in-sample,
$\valAUCFLfirmTemp$ under temporal holdout) and direct-defendant
discrimination ($\valAUCdirectStd$) is the design's empirical
signature, not a failure of validation. The screening statistic
recovers loser-side participation footprints; direct defendants
are by construction the winner side of the same arrangements
and lie outside the target population.

\subsection{Robustness to the enforcement record}
\label{sec:cade_robustness}

Dropping all $\valCADEpermPerm$ CADE-involved tender-items and
re-estimating the within-PBU baseline yields $\hat\beta = 0.062$
($N = \valCADEpermObs{,}954$), virtually unchanged from
$\valPctSixfour$. The screening signal in $\beta$ operates
independently of the enforcement record. The residual sample
either contains undetected cartels that the screening statistic
flags before enforcement reaches them, or higher-price
environments correlated with frequent-loser presence through
channels orthogonal to documented collusion. Both readings are
consistent with the screening interpretation, and neither
relies on the within-sample CADE adjudications for its content.

\subsection{What the validation does and does not show}
\label{sec:cade_scope}

The validation supports four claims. First, the screening
statistic concentrates on cobidder-labeled firms inside the
always-loser stratum at a $\valMultThreeTwo\!\!\times$
contemporaneous excess (AUC $\valAUCprePost$) and a
$\valMultThreeFive\!\!\times$ prospective excess (AUC
$\valAUCFLfirm$ in-sample, $\valAUCFLfirmTemp$ under temporal
holdout). Second, the price imprint $\beta$ is insensitive to
removing CADE-involved tender-items, which means the screening
signal is not driven by the enforcement record itself. Third,
the asymmetry between cobidder and direct-defendant
discrimination delimits the scope: the screening statistic
recovers cartel-adjacent environments, not cartel members.
Fourth, the labels are indirect by construction---a feature of
working at the award layer rather than the bid layer. The
validation corroborates the screening interpretation of the
construct; it does not, and under coarsened observability cannot,
adjudicate the legal status of any individual firm.
